\section{Problem Setup}

We study minimum-time trajectory optimization for a vehicle on a closed track with optional static circular obstacles. The optimizer uses a spatial discretization with horizon length $N=120$ on the Oval track. The underlying state used by the optimizer is
\[
x = [u_x, u_y, r, \Delta F_{z,\mathrm{long}}, \Delta F_{z,\mathrm{lat}}, t, e, d\psi],
\]
with controls
\[
u = [\delta, F_x].
\]

The optimization objective is the track lap-time objective plus a small control regularizer. Track bounds and obstacle-clearance constraints are enforced along the spatial trajectory. The DT does not replace this optimizer. Instead, it predicts an initial state-control rollout that is either accepted and passed to the optimizer or rejected and replaced by the baseline initializer.

The core evaluation question is therefore:
\begin{quote}
Does the DT warm-start reduce downstream solver work relative to the baseline initializer, while preserving solve success?
\end{quote}
